\chapter{Literature Review}
\label{ch:litreview}

This chapter surveys the academic literature relevant to \projname{}.
It is organised by theme: (i)~AutoML systems, (ii)~hyperparameter
optimisation and neural architecture search, (iii)~gradient-boosted
trees on tabular data, (iv)~model explainability, (v)~LLM tool use
and agentic systems, (vi)~the Model Context Protocol, and
(vii)~conversational data-science interfaces. Only verifiable
publications are cited; every reference has a corresponding entry in
\file{bibliography/references.bib}.

\section{AutoML Systems}
\label{sec:lit-automl}

Auto-WEKA \citep{thornton2013autoweka} was among the first systems to
frame AutoML as a joint choice of algorithm and hyperparameters
solved by Bayesian optimisation. Auto-sklearn
\citep{feurer2015autosklearn} extended this by meta-learning warm
starts and constructing ensembles. TPOT \citep{olson2016tpot} instead
uses genetic programming over scikit-learn pipelines. H2O AutoML
\citep{ledell2020h2o} is a production-oriented Java/Python system
that combines random search over gradient-boosted trees and
generalised linear models with stacked ensembles. AutoGluon
\citep{erickson2020autogluon} argues that bagging and stacking of
heterogeneous models are more effective on tabular data than
NAS-style search. Auto-Keras \citep{jin2019autokeras} focuses on
architecture search for deep image models.

More recent contributions include CAAFE \citep{hollmann2024caafe},
which uses an LLM to propose new features in an interpretable
feature-engineering loop, and DS-Agent \citep{guo2024dsagents}, which
uses case-based reasoning to plan data-science workflows. \projname{}
inherits from this fifth-generation line but focuses on
\emph{platform} concerns --- deployment, artefact management, and
extensibility --- rather than novel search algorithms.

\section{Hyperparameter Optimisation and NAS}
\label{sec:lit-hpo}

\citet{bergstra2012random} demonstrated that random search often
outperforms grid search on high-dimensional configuration spaces.
SMAC \citep{hutter2011smac} popularised sequential model-based
optimisation with random-forest surrogates. TPE
\citep{bergstra2011algorithms} models the density of good and bad
configurations separately and forms the default sampler in Optuna
\citep{akiba2019optuna}. Hyperband \citep{li2017hyperband} and BOHB
\citep{falkner2018bohb} use multi-fidelity budgets to accelerate
search on expensive learners. Neural architecture search
\citep{zoph2017nas, real2019regularized, elsken2019nas} operates in
much larger, discrete search spaces than tabular HPO and is out of
scope for the current \projname{} implementation.

\section{Gradient-Boosted Trees on Tabular Data}
\label{sec:lit-tabular}

Gradient boosting \citep{friedman2001greedy} remains the dominant
family of methods on tabular data. XGBoost \citep{chen2016xgboost}
introduced sparsity-aware split finding and regularisation that
enabled scaling to millions of rows. LightGBM \citep{ke2017lightgbm}
proposed histogram binning, gradient-based one-side sampling, and
exclusive-feature bundling to reduce training time. CatBoost
\citep{prokhorenkova2018catboost} handled categorical features with
target-based encodings. \citet{grinsztajn2022trees} conducted a
comprehensive empirical comparison and concluded that tree ensembles
continue to outperform deep models on medium-sized tabular data;
\citet{borisov2022deep} report similar findings. Random forests
\citep{breiman2001rf} remain a strong, hyperparameter-tolerant
baseline and are always included in the \projname{} bake-off.

\section{Explainability}
\label{sec:lit-xai}

\citet{ribeiro2016lime} proposed LIME, a local surrogate
explanation. \citet{lundberg2017shap} introduced SHAP as a unified
framework that satisfies desirable game-theoretic axioms and
\citet{lundberg2020treeshap} presented an $O(TL D^{2})$ algorithm
for exact Shapley values on tree ensembles. \projname{} uses
\code{TreeExplainer} for tree-based champions and falls back to
\code{KernelExplainer} for linear and neural models. For a survey
of interpretable ML, see \citet{molnar2022ibm}. The distinction
between explanation and interpretation is critically discussed by
\citet{rudin2019stop}.

\section{Data Validation and MLOps}
\label{sec:lit-mlops}

Schema validation for pandas dataframes is provided by pandera
\citep{pandera}, which allows column-level type, nullability, and
categorical constraints. \citet{polyzotis2018dataval,
schelter2018data} argue that data validation is a critical
MLOps concern that is under-served by the sklearn ecosystem.
Reproducibility guidelines for ML papers are surveyed by
\citet{pineau2021reproducibility}. Feature stores, experiment
trackers, and model registries have become standard concerns in
production ML but are largely out of scope for a single-user
research platform.

\section{LLM Tool Use, Function Calling, and Agents}
\label{sec:lit-llm-tools}

Instruction-tuned language models such as GPT-3
\citep{brown2020gpt3}, GPT-4 \citep{openai2023gpt4}, and Llama\,3
\citep{llama3card, touvron2023llama} exhibit strong zero-shot
capabilities that are further improved by tool use. Toolformer
\citep{schick2023toolformer} showed that a model can be
self-supervised to invoke external APIs. ReAct \citep{yao2023react}
interleaves reasoning traces with tool calls to improve
compositionality. Gorilla \citep{patil2023gorilla} explores
retrieval-augmented tool selection over large API catalogs. OpenAI's
function-calling API \citep{openai2023function} standardised the
tool-call schema that is used by Groq and other OpenAI-compatible
providers. \citet{qin2024toolllm} and \citet{shen2024hugginggpt}
extend agentic tool use to thousands of external services.

\projname{} is deliberately more conservative than these systems:
its tool catalog is a curated set of eleven functions and its
planner is a single LLM without a hierarchical controller or a
memory module. This restricts the surface area for hallucination and
keeps the platform inspectable.

\section{The Model Context Protocol}
\label{sec:lit-mcp}

MCP was announced in late 2024 \citep{anthropic2024mcp} and
formalised in an open specification \citep{mcpSpec2024}. The
specification defines the JSON-RPC-based transport, the three
primitive kinds (tools, resources, prompts), the discovery lifecycle,
and the sampling protocol. FastMCP \citep{fastmcp} is a
community-maintained Python implementation that produces MCP
schemas from decorated Python functions.

Academic study of MCP is still nascent at the time of writing, but
the protocol is closely related to earlier work on tool exchange
formats \citep{schick2023toolformer, patil2023gorilla} and to the
OpenAI function-calling schema \citep{openai2023function}. The
particular contribution of MCP is that it decouples the
tool-provider from the LLM host --- a property that matches the
service-oriented architecture used by \projname.

\section{Conversational Interfaces for Data Science}
\label{sec:lit-convo}

The idea of a chat-first interface to data analysis predates modern
LLMs; work such as Datatone \citep{gao2015datatone} and
\citet{narechania2020nl4dv} used semantic parsing to translate
natural-language questions into Vega-Lite specifications
\citep{satyanarayan2016vega}. LIDA \citep{dibia2023lida} and Chat2Vis
\citep{maddigan2023chat2vis} use LLMs to synthesise visualisation
code. LLM-based ``code interpreter'' environments (as popularised by
ChatGPT) let the model both generate and execute code in a sandbox
but do not expose a typed tool catalog to a third-party developer.
\projname{} differs from these systems by exposing an
\emph{explicit, extensible, and well-typed} catalog that is versioned
inside the repository rather than embedded in a proprietary sandbox.

\section{Positioning of \projname}
\label{sec:lit-position}

Combining the above threads, \projname{} occupies the following
position in the design space:

\begin{itemize}[leftmargin=1.4em]
  \item Like Auto-sklearn \citep{feurer2015autosklearn} and AutoGluon
        \citep{erickson2020autogluon}, it uses cross-validated model
        selection with the standard tabular estimator family.
  \item Like Optuna \citep{akiba2019optuna}, it exposes budget-bounded
        TPE search.
  \item Like SHAP \citep{lundberg2017shap}, it provides model-agnostic
        explanations of the champion model.
  \item Unlike Auto-sklearn or AutoGluon, its \emph{interface} is a
        conversation and its \emph{orchestrator} is an LLM that
        selects and composes tools.
  \item Unlike code-interpreter systems, its tool catalog is a
        versioned, typed set of MCP tools that can be inspected,
        replaced, and extended without touching the front-end.
  \item Unlike commercial AutoML SaaS, all storage is on the local
        filesystem and no authentication or account is required.
\end{itemize}

The remainder of the report elaborates on this design.
